%%%%%%%%%%%%%%%%%%%%%%%%%%%%%%%%%%%%%%%%%%%%%%%%%%%%%%%%%%%%%%%%%%%%%%%%%%%%%%%%
% resumo.tex
%
% Modelo de arquivo para uso com a classe uvvTeX2, para a formatação de
% trabalhos acadêmicos na Universidade Vila Velha (UVV) (https://www.uvv.br).
%
% Para maiores informações, visite:
%    https://github.com/uvv-computacao/uvvtex2
%
% Neste arquivo você deve escrever o resumo de sua monografia, ou seja, deve
% escrever o resumo em PORTUGUÊS. O resumo deve ser objetivo e listar os
% aspectos mais importantes de seu trablaho. Ao final do abstract, as
% palavras-chave que você definiu no arquivo principal da monografia serão
% inseridas automaticamente.
%%%%%%%%%%%%%%%%%%%%%%%%%%%%%%%%%%%%%%%%%%%%%%%%%%%%%%%%%%%%%%%%%%%%%%%%%%%%%%%%

% Não altere as linhas a seguir:
\setlength{\absparsep}{18pt}
\begin{resumo}

% Começe a escrever o abstract aqui:
Este trabalho apresenta a análise e a modelagem de uma solução de automação do
fluxo pedagógico no sistema Educa EMES, utilizado pela Escola da Magistratura do
Espírito Santo (EMES). O estudo parte da existência de uma plataforma integrada
já em produção, mas identifica lacunas relacionadas à permanência de
intervenções manuais em etapas como abertura e encerramento de inscrições,
confirmação de participantes, aprovação, certificação e geração de relatórios.
Metodologicamente, a pesquisa caracteriza-se como aplicada, de abordagem
qualitativa e baseada em estudo de caso, tendo como objeto um sistema real
utilizado em contexto institucional. São analisados o fluxo pedagógico atual, a
estrutura tecnológica e o modelo de dados do sistema, os principais gargalos
operacionais e os requisitos necessários para ampliar a automação. Como
resultado, propõe-se uma modelagem baseada em etapas
configuráveis, análise de elegibilidade, identificação de pendências,
rastreabilidade dos estados e preservação da possibilidade de intervenção
manual, orientando as etapas de desenvolvimento e avaliação da solução.

% Não altere as linhas a seguir:
\textbf{Palavras-chave}: \imprimirpalavraschave.
\end{resumo}
